\chapter{Numerical solution}
\label{cha:numerical}

This chapter describes how the continuous model of Chapter~\ref{cha:formulation}
is turned into an algorithm: the data structures, the explicit time integrator
and its stability limit, the contact-detection strategy, the bookkeeping of the
tangential history, the deposition procedure that builds the initial packing,
and the resulting computational cost.

\section{Overview of the algorithm}
\label{sec:overview}

The simulation is advanced by a fixed-step explicit scheme.  One time step
consists of the operations listed below; the numbering is used as a reference in
the remainder of the chapter.

\begin{center}
\fbox{\parbox{0.93\textwidth}{\small
\textbf{Algorithm 1 -- one time step} (\texttt{DEMSimulation.step})
\begin{enumerate}[leftmargin=2.2em,nosep]
  \item \emph{First half kick} (uses the forces stored from the previous
        evaluation)\\
        $\mathbf{v}_{i}\!\leftarrow\!\mathbf{v}_{i}
          +\tfrac{\Delta t}{2}\mathbf{F}_{i}/m_{i}$,\quad
        $\omega_{i}\!\leftarrow\!\omega_{i}+\tfrac{\Delta t}{2}T_{i}/I_{i}$,\quad
        $v_{p}\!\leftarrow\!v_{p}+\tfrac{\Delta t}{2}a_{p}\bigl(t^{n}\bigr)$
  \item \emph{Drift}\\
        $\mathbf{x}_{i}\!\leftarrow\!\mathbf{x}_{i}+\Delta t\,\mathbf{v}_{i}$,\quad
        $y_{p}\!\leftarrow\!y_{p}+\Delta t\,v_{p}$,\quad
        $t\!\leftarrow\!t+\Delta t$
  \item \emph{Force evaluation} (\texttt{compute\_forces})
    \begin{enumerate}[leftmargin=1.6em,nosep,label=(\alph*)]
      \item reset $\mathbf{F}_{i}$, $T_{i}$; add gravity $-m_{i}g\,\mathbf{e}_{y}$;
      \item rebuild the cell list and generate candidate pairs
            (Section~\ref{sec:broadphase});
      \item for each candidate pair: exact overlap test~\eqref{eq:overlap},
            contact law of Section~\ref{sec:contact-summary}, assembly of
            $\pm\mathbf{F}_{ij}$ and $T_{ij},T_{ji}$;
      \item for each particle: three wall contacts and, if active, the platen
            contact; accumulate the platen reaction~\eqref{eq:platen-reaction};
      \item replace the tangential-history table by the one just built
            (Section~\ref{sec:history}).
    \end{enumerate}
  \item \emph{Second half kick} (uses the new forces and $F(t^{n+1})$)\\
        $\mathbf{v}_{i}\!\leftarrow\!\mathbf{v}_{i}
          +\tfrac{\Delta t}{2}\mathbf{F}_{i}/m_{i}$,\quad
        $\omega_{i}\!\leftarrow\!\omega_{i}+\tfrac{\Delta t}{2}T_{i}/I_{i}$,\quad
        $v_{p}\!\leftarrow\!v_{p}+\tfrac{\Delta t}{2}a_{p}\bigl(t^{n+1}\bigr)$
\end{enumerate}
}}
\end{center}

Three drivers chain these steps into the physical stages of the experiment:
\texttt{rain\_until\_full} (deposition, Section~\ref{sec:deposition}),
\texttt{rest} (relaxation) and \texttt{compress} (loading).

\section{State representation}
\label{sec:state}

All particle data are held in contiguous \texttt{NumPy} arrays that grow as
particles are inserted:

\begin{center}
\begin{tabular}{lll}
  \toprule
  Array & Shape & Content \\
  \midrule
  \texttt{pos}, \texttt{vel}, \texttt{force} & $(N,2)$ & $\mathbf{x}_{i}$, $\mathbf{v}_{i}$, $\mathbf{F}_{i}$ \\
  \texttt{omega}, \texttt{torque}            & $(N,)$  & $\omega_{i}$, $T_{i}$ \\
  \texttt{radius}, \texttt{mass}, \texttt{inertia} & $(N,)$ & $R_{i}$, $m_{i}$, $I_{i}$ \\
  \texttt{inv\_mass}, \texttt{inv\_inertia}  & $(N,)$  & $1/m_{i}$, $1/I_{i}$ \\
  \texttt{tangential}                        & dict    & $\boldsymbol{\xi}$ per contact \\
  \bottomrule
\end{tabular}
\end{center}

Two implementation details matter.  First, the reciprocals $1/m_{i}$ and
$1/I_{i}$ are cached because they are used twice per step, in the two half
kicks.  Second, new particles are always \emph{appended} at the end of the
arrays, never inserted; existing indices are therefore stable, which is what
makes the index-keyed history table of Section~\ref{sec:history} valid across
insertions.  Whenever the radius array changes, $m_{i}$, $I_{i}$ and their
reciprocals are recomputed from~\eqref{eq:massprops}.

\section{Time integration}
\label{sec:integration}

\subsection{Velocity-Verlet scheme}

The equations of motion~\eqref{eq:newton}--\eqref{eq:euler} are integrated with
the velocity-Verlet algorithm~\cite{swope1982,allen2017} in its kick--drift--kick
form:
\begin{align}
  \mathbf{v}^{\,n+1/2} &= \mathbf{v}^{\,n}
    + \frac{\Delta t}{2}\,\frac{\mathbf{F}^{\,n}}{m},
    \label{eq:vv1}\\
  \mathbf{x}^{\,n+1} &= \mathbf{x}^{\,n} + \Delta t\,\mathbf{v}^{\,n+1/2},
    \label{eq:vv2}\\
  \mathbf{F}^{\,n+1} &= \mathbf{F}\bigl(\mathbf{x}^{\,n+1},
                        \mathbf{v}^{\,n+1/2},
                        \boldsymbol{\xi}^{\,n+1}\bigr),
    \label{eq:vv3}\\
  \mathbf{v}^{\,n+1} &= \mathbf{v}^{\,n+1/2}
    + \frac{\Delta t}{2}\,\frac{\mathbf{F}^{\,n+1}}{m}.
    \label{eq:vv4}
\end{align}
The rotational degrees of freedom and the platen degree of freedom are advanced
by the same three-stage pattern, with $T_{i}/I_{i}$ and
$a_{p}=\bigl(F_{p}-F(t)-m_{p}g_{p}\bigr)/m_{p}$ in place of $\mathbf{F}/m$.

The scheme is explicit, requires a single force evaluation per step, stores no
history beyond the current state, is time reversible, and is second-order
accurate: the local truncation error is $O(\Delta t^{3})$ in position and the
global error is $O(\Delta t^{2})$.  For conservative forces it is symplectic,
so the energy error remains bounded rather than drifting -- an essential
property for the long deposition runs.

\subsection{Velocity-dependent forces}

The contact force is not conservative: it depends on $\mathbf{v}$ through the
dashpots~\eqref{eq:cn} and~\eqref{eq:Ft-trial}.  In~\eqref{eq:vv3} it is
evaluated at the half-step velocity $\mathbf{v}^{\,n+1/2}$, which is the only
velocity available at that point without an implicit solve or an extra
iteration.  Since
$\mathbf{v}^{\,n+1/2}=\mathbf{v}^{\,n+1}-\tfrac{\Delta t}{2}\mathbf{a}^{\,n+1}
 +O(\Delta t^{2})$, the dissipative part of the force carries an $O(\Delta t)$
error, so the scheme is formally first order in the damping contribution and
strictly symplectic only in the limit $\beta\to0$.  The resulting bias is
proportional to $\beta$ and to $\Delta t/t_{c}$; with the values used here
($\beta=0.28$, $\Delta t/t_{c}\approx0.03$) it is well below the other
modelling errors.  The same approximation is made by most production DEM
codes~\cite{silbert2001,kloss2012}.

\subsection{Reference integrator for verification}

The class \texttt{VelocityVerlet} implements~\eqref{eq:vv1}--\eqref{eq:vv4} for
a generic acceleration callback $\mathbf{a}(\mathbf{x},\mathbf{v})$.  It is not
used by the production loop -- which fuses the kicks with the force evaluation
to avoid recomputing $\mathbf{F}^{\,n}$ -- but it exposes the integrator to unit
testing against analytical solutions (Section~\ref{sec:unit-tests}).

\section{Stable time step}
\label{sec:timestep}

Being explicit, the integrator is only conditionally stable.  The binding
restriction is the fastest elastic response of the assembly, which is the
propagation of a surface (Rayleigh) wave across the \emph{smallest} particle.
The Rayleigh critical time step is
\begin{equation}
  \boxed{\;
  T_{R} = \frac{\pi R_{\min}}{0.1631\,\nu + 0.8766}
          \sqrt{\frac{\rho}{G}} \;}
  \label{eq:rayleigh}
\end{equation}
where the denominator is the standard rational approximation of the ratio
between the Rayleigh and the shear wave speeds~\cite{osullivan2004}.  The code
uses a safety fraction
\begin{equation}
  \Delta t = \alpha\,T_{R},
  \qquad \alpha = 0.2 ,
  \label{eq:dt}
\end{equation}
which lies in the middle of the range $0.1$--$0.3$ recommended in the
literature for dense assemblies with coordination numbers of four to six.  A
user-supplied $\Delta t$ overrides~\eqref{eq:dt}.

Three remarks on~\eqref{eq:rayleigh}:
\begin{itemize}[leftmargin=1.6em]
  \item It scales with $R_{\min}$, not with the mean radius.  A single very
        small grain therefore penalises the whole simulation, which is why the
        size ratio $R_{\max}/R_{\min}$ should be kept moderate.
  \item It scales with $\sqrt{\rho/G}$, i.e.\ with the inverse shear-wave
        speed.  Reducing $E$ (and hence $G$) by a factor of $100$ increases the
        admissible time step by a factor of $10$; this "artificial softening" is
        the standard device that makes DEM of stiff materials affordable, at
        the cost of larger overlaps (Section~\ref{sec:limitations}).
  \item It is independent of the load.  A complementary check is that
        $\Delta t$ must resolve the collision time~\eqref{eq:tc}; both criteria
        are evaluated in Table~\ref{tab:dt} and the Rayleigh one is the more
        restrictive.
\end{itemize}

\begin{table}[htbp]
  \centering
  \caption{Time-step data for the reference simulation
  ($R_{\min}=\qty{4}{\milli\metre}$, $R_{\max}=\qty{20}{\milli\metre}$).}
  \label{tab:dt}
  \begin{tabular}{llS[table-format=1.3e2]l}
    \toprule
    Symbol & Quantity & {Value} & Unit \\
    \midrule
    $T_{R}$    & Rayleigh critical step, Eq.~\eqref{eq:rayleigh} & 4.843e-4 & \unit{\second} \\
    $\alpha$   & safety fraction                                 & 2.000e-1 & -- \\
    $\Delta t$ & time step used                                  & 9.686e-5 & \unit{\second} \\
    $t_{c}$    & collision time at $v_{0}=\qty{1}{\metre\per\second}$,
                 Eq.~\eqref{eq:tc}                               & 3.269e-3 & \unit{\second} \\
    $t_{c}/\Delta t$ & steps resolving one collision             & 3.380e1  & -- \\
    \bottomrule
  \end{tabular}
\end{table}

\section{Contact detection}
\label{sec:broadphase}

\subsection{Uniform grid (cell list)}

A naive search over all $\tfrac{1}{2}N(N-1)$ pairs costs $O(N^{2})$ per step and
is unaffordable.  Instead the domain is covered by a uniform grid of square
cells of side
\begin{equation}
  h = 2R_{\max},
  \label{eq:cellsize}
\end{equation}
and every particle is binned into the cell
\begin{equation}
  c(\mathbf{x}_{i}) = \bigl(\lfloor x_{i}/h\rfloor,\;
                            \lfloor y_{i}/h\rfloor\bigr) \in \mathbb{Z}^{2}.
  \label{eq:cellindex}
\end{equation}
Only occupied cells are stored, in a hash map from the index pair to the list of
particles it contains, so the memory cost is $O(N)$ regardless of how empty the
box is -- an important property during deposition, when the bed occupies a small
and growing fraction of the domain.

\paragraph{Correctness.}  The choice~\eqref{eq:cellsize} guarantees that no
contact is missed.  If $i$ and $j$ touch then
$d_{ij}<R_{i}+R_{j}\le 2R_{\max}=h$, hence
$\lvert x_{i}-x_{j}\rvert<h$ and $\lvert y_{i}-y_{j}\rvert<h$, so their cell
indices differ by at most one in each direction.  Searching a particle's own
cell and its eight neighbours is therefore exhaustive.

\paragraph{Half stencil.}  To visit every unordered pair exactly once, the
search combines the pairs inside each cell with the pairs between the cell and
only \emph{half} of its neighbours,
\begin{equation}
  \mathcal{O} = \bigl\{(1,0),\ (0,1),\ (1,1),\ (-1,1)\bigr\},
  \label{eq:stencil}
\end{equation}
the remaining four offsets being the negatives of these and therefore already
covered when the neighbouring cell is processed.  This halves the number of
candidate pairs and, more importantly, removes the risk of applying a contact
force twice.

\paragraph{Narrow phase.}  Each candidate pair is then subjected to the exact
test~\eqref{eq:overlap}; non-overlapping candidates are discarded before any
constitutive evaluation.  The pair is canonically ordered as $(i,j)$ with $i<j$
so that the history key of Section~\ref{sec:history} is unique.

\subsection{Cost and the effect of polydispersity}

With a cell size fixed by the largest particle, the expected number of particles
per cell is
\begin{equation}
  n_{\text{cell}} \sim \phi \left(\frac{h}{2\bar{R}}\right)^{2}
  \quad\text{and, in the worst case,}\quad
  n_{\text{cell}}^{\max} \sim
  \left(\frac{R_{\max}}{R_{\min}}\right)^{2},
\end{equation}
so the number of candidate pairs per particle grows quadratically with the size
ratio.  For the reference case $R_{\max}/R_{\min}=5$, which is already a
noticeable penalty: a hierarchical grid or per-size-class cell lists would be
the natural remedy for broader distributions.  The overall cost of one step is
nevertheless
\begin{equation}
  \mathcal{O}\bigl(N\bigr)\ \text{(binning)}
  + \mathcal{O}\bigl(N\,\bar{n}_{\text{pairs}}\bigr)\ \text{(narrow phase)},
\end{equation}
i.e.\ linear in $N$ with a size-ratio-dependent constant.  The grid is rebuilt
from scratch at every step; no Verlet neighbour list with a skin distance is
used, which keeps the implementation simple and free of the "lost contact"
errors that an out-of-date neighbour list can produce.

\section{Tangential-history bookkeeping}
\label{sec:history}

The tangential spring $\boldsymbol{\xi}$ must persist as long as the contact
exists and must disappear with it.  This is achieved with a dictionary keyed by
a contact identifier:
\begin{center}
\begin{tabular}{ll}
  \toprule
  Key & Contact \\
  \midrule
  \texttt{("pp", i, j)} & particle $i$ -- particle $j$, with $i<j$ \\
  \texttt{("wL", i)}, \texttt{("wR", i)}, \texttt{("wB", i)} & particle $i$ -- left / right / bottom wall \\
  \texttt{("pl", i)} & particle $i$ -- platen \\
  \bottomrule
\end{tabular}
\end{center}

At every force evaluation a \emph{fresh} table is built and the old one is
consulted through
$\boldsymbol{\xi}_{\text{old}}=\texttt{old.get(key},\ \mathbf{0})$.  This single
design decision provides three behaviours for free:

\begin{enumerate}[leftmargin=2.2em]
  \item \textbf{Creation.}  A contact that closes for the first time finds no
        entry and starts from $\boldsymbol{\xi}=\mathbf{0}$, as required.
  \item \textbf{Persistence.}  A contact that survives copies its spring
        forward, which is projected and incremented
        by~\eqref{eq:xi-project}--\eqref{eq:xi-accumulate}.
  \item \textbf{Deletion.}  A contact that opens is simply not re-inserted into
        the new table; its history is garbage collected, and a later re-closure
        starts afresh.  No explicit pruning pass is needed, and the memory
        occupied by the table is always proportional to the number of
        \emph{active} contacts.
\end{enumerate}

The keys are index based, which is safe only because indices are never
permuted (Section~\ref{sec:state}).  The price of the approach is one dictionary
allocation per step and the cost of hashing every active contact, which is
significant in Python but not in the overall $O(N)$ picture.

Two consistency requirements follow from~\eqref{eq:xi-accumulate}:
\texttt{compute\_forces} must be called exactly once per time step, and the
$\Delta t$ passed to the contact law must be the integration step.  The drivers
call \texttt{compute\_forces} once before entering their loop to prime
$\mathbf{F}^{\,0}$ for the first half kick; this initial evaluation contributes
one extra increment to the springs of the contacts that are already closed, an
$O(\Delta t)$ perturbation of a quiescent state that is immediately relaxed.

\section{Boundary and platen coupling}
\label{sec:coupling}

Wall contacts are evaluated in a single pass over the particles, using the
overlaps of Table~\ref{tab:boundaries} and the degenerate effective quantities
\eqref{eq:wall-effective}.  Because the three fixed walls are tested for every
particle irrespective of its position, the cost is $O(N)$ with a very small
constant; the overlap test rejects the interior particles immediately.

The platen is integrated with the same velocity-Verlet pattern as the particles,
in a \emph{staggered} (explicit, one-way at each half step) fashion:
\begin{itemize}[leftmargin=1.6em]
  \item the first half kick uses the reaction $F_{p}$ computed at the end of the
        previous step and the load $F(t^{n})$;
  \item the drift moves the platen with $v_{p}^{\,n+1/2}$, so the particles see
        the new platen position and the half-step platen velocity
        in~\eqref{eq:vrel-platen};
  \item the force evaluation accumulates the new reaction
        $F_{p}^{\,n+1}=-\sum_{i}\mathbf{F}_{ip}\cdot\mathbf{e}_{y}$;
  \item the second half kick uses $F_{p}^{\,n+1}$ and $F(t^{n+1})$.
\end{itemize}
The coupling is therefore fully explicit and adds no stability restriction of
its own provided that the platen is not much lighter than the grains it
touches; a very light platen would behave like an extra particle with a very
small mass and would require a smaller $\Delta t$ through the $m^{*}$
in~\eqref{eq:cn}.  A practical rule is $m_{p}\gtrsim 10\,m_{\max}$; the
reference case uses $m_{p}=\qty{0.5}{\kilogram}$ against
$m_{\max}=\qty{0.084}{\kilogram}$.

The load clock \texttt{time\_load} is advanced only while the platen is active,
so $F(t)$ in~\eqref{eq:ramp} always starts from zero when compression begins,
independently of how long the deposition took.

\section{Initial configuration and deposition}
\label{sec:deposition}

A DEM result is only as good as its initial packing.  Two generators are
provided.

\subsection{Loose lattice generator}

\texttt{generate\_rain\_particles} places $N$ particles on a jittered
rectangular lattice above the bed.  With a lattice pitch
\begin{equation}
  s = 2 R_{\max}\,\sigma, \qquad \sigma > 1
  \quad(\text{default } \sigma = 1.15),
\end{equation}
the number of columns that fit between the side walls is
$n_{\text{col}}=\lfloor (B-2R_{\max})/s \rfloor + 1$, and particle $k$ is
assigned to column $k \bmod n_{\text{col}}$ and row $\lfloor k/n_{\text{col}}
\rfloor$.  A uniform random offset of amplitude
\begin{equation}
  \Delta_{\text{jitter}} = \tfrac{1}{2}\,j\,\bigl(s - 2R_{\max}\bigr),
  \qquad 0 \le j \le 1 \quad (\text{default } j = 0.25),
\end{equation}
is added to both coordinates.  The bound on the amplitude is chosen so that
\emph{no} initial overlap can be created, whatever the random draw: the residual
clearance $s-2R_{\max}$ is never fully consumed.  Breaking the lattice symmetry
matters, because a perfect lattice collapses into a crystalline packing with
spurious anisotropy.  Radii are drawn from~\eqref{eq:radii} with an explicit
seed, so packings are reproducible.

\subsection{Wave-by-wave raining}

The driver actually used for the reference simulation,
\texttt{rain\_until\_full}, starts from an \emph{empty} box and fills it
adaptively.

\begin{center}
\fbox{\parbox{0.93\textwidth}{\small
\textbf{Algorithm 2 -- raining until full} (\texttt{rain\_until\_full})
\begin{enumerate}[leftmargin=2.2em,nosep]
  \item $w \leftarrow 0$ (counter of consecutive empty waves)
  \item \textbf{while} $w < w_{\max}$ \textbf{and} $N < N_{\max}$
        \textbf{and} steps $<$ cap:
    \begin{enumerate}[leftmargin=1.6em,nosep,label=(\alph*)]
      \item attempt to insert $n_{b}$ particles in the band
            $y\in[L-2R_{\max},\,L-R_{\max}]$, $x\in[R,\,B-R]$, by rejection
            sampling: a candidate is accepted only if it overlaps neither an
            existing particle near the top nor a particle already accepted in
            this wave; up to $n_{a}$ attempts per particle;
      \item append the accepted particles with velocity
            $(0,-v_{\text{rain}})$;
      \item $w \leftarrow 0$ if at least one particle was inserted, else
            $w \leftarrow w+1$;
      \item advance $n_{s}$ time steps (Algorithm 1).
    \end{enumerate}
  \item \textbf{final settling}: advance up to $n_{\text{settle}}$ steps,
        stopping as soon as the quiescence criterion~\eqref{eq:quiet} is met.
\end{enumerate}
}}
\end{center}

The rejection test is restricted to the particles whose centres lie above
$L-4R_{\max}$, which keeps the insertion cost proportional to the number of
grains near the surface rather than to $N$.  The small downward seed velocity
$v_{\text{rain}}$ removes the initial free-fall transient and makes the arrival
of successive waves more regular.  The termination logic is purely geometric:
when the growing bed reaches the insertion band, rejection sampling can no
longer place anything, $w$ reaches $w_{\max}$ and the loop exits.  The number of
particles is therefore an \emph{output} of the procedure, determined by the box
geometry and the size distribution, not an input.  Default values are collected
in Table~\ref{tab:deposition}.

\begin{table}[htbp]
  \centering
  \caption{Deposition parameters (values used in the reference simulation).}
  \label{tab:deposition}
  \begin{tabular}{llll}
    \toprule
    Symbol & Meaning & Code name & Value \\
    \midrule
    $n_{b}$            & particles per wave                & \texttt{batch\_size}     & 15 \\
    $n_{s}$            & steps between waves               & \texttt{steps\_between}  & 450 \\
    $n_{a}$            & placement attempts per particle   & \texttt{attempts}        & 30 \\
    $w_{\max}$         & empty waves before stopping       & \texttt{patience}        & 4 \\
    $v_{\text{rain}}$  & seed velocity                     & \texttt{\_rain\_speed}   & \qty{0.3}{\metre\per\second} \\
    $n_{\text{settle}}$& maximum settling steps            & \texttt{settle\_steps}   & 4000 \\
    $N_{\max}$         & hard cap on particle number       & \texttt{max\_particles}  & 2000 \\
    $\epsilon$         & quiescence tolerance              & \texttt{quiet\_tol}      & \qty{1e-6}{\joule} \\
    -- & random seed & \texttt{seed} & 7 \\
    \bottomrule
  \end{tabular}
\end{table}

\subsection{Quiescence criterion}

Deposition and settling terminate on the mean translational kinetic energy per
particle,
\begin{equation}
  \boxed{\;
  \bar{E}_{k}
  = \frac{1}{N}\sum_{i=1}^{N}\tfrac{1}{2}m_{i}\lVert\mathbf{v}_{i}\rVert^{2}
  \;<\; \epsilon \;}
  \label{eq:quiet}
\end{equation}
combined with a minimum number of steps, so that the criterion cannot be
satisfied trivially at $t=0$ when every particle is still at rest.  The
normalisation by $N$ makes the threshold independent of how many particles have
been deposited, which is essential here because $N$ grows during the run.  A
useful way to choose $\epsilon$ is to compare it with the potential energy of a
grain displaced by a small fraction of its own radius: with
$\epsilon=\qty{1e-6}{\joule}$ and a mean grain mass of about
\qty{2e-2}{\kilogram}, the corresponding root-mean-square velocity is
$\sqrt{2\epsilon/\bar{m}}\approx\qty{1e-2}{\metre\per\second}$, i.e.\ two orders
of magnitude below the raining velocity.

\section{Recording and post-processing}
\label{sec:recording}

Two output streams are collected at user-defined intervals.

\begin{description}[leftmargin=1.6em,style=nextline]
  \item[Frames] $\bigl(t,\ \mathbf{x}_{i},\ R_{i},\
        \lVert\mathbf{v}_{i}\rVert,\ y_{p},\ \text{phase}\bigr)$, recorded every
        \texttt{record\_every} steps.  Positions and radii are \emph{copied},
        not referenced, because the underlying arrays are overwritten in place
        and reallocated when particles are appended.  The phase label
        (\texttt{deposition}, \texttt{settling}, \texttt{loading}) is carried
        along so that the animation can be annotated.
  \item[Load curve] $\bigl(t,\ F(t),\ u_{p},\ F_{p}\bigr)$, recorded only during
        the compression stage.  Comparing $F(t)$ with $F_{p}$ is the simplest
        available check of quasi-staticity: the two curves coincide when the
        platen inertia is negligible, and separate when the loading is too fast.
\end{description}

The number of stored frames, not the number of time steps, dominates the memory
footprint of a long run: each frame holds $O(N)$ floating-point numbers.

\section{Verification strategy}
\label{sec:verification}

The implementation is verified at three levels, in increasing order of
integration:
\begin{enumerate}[leftmargin=2.2em]
  \item \emph{Unit level}: each closed-form expression
        (\eqref{eq:moduli}, \eqref{eq:hertz}, \eqref{eq:coulomb}) and the
        integrator are compared with analytical results
        (Section~\ref{sec:unit-tests});
  \item \emph{Component level}: single-contact experiments -- a bouncing grain,
        a sliding grain -- reproduce the prescribed restitution and friction
        coefficients (Section~\ref{sec:component-tests});
  \item \emph{System level}: the full deposition-and-compression experiment is
        checked for global consistency, in particular the equilibrium condition
        $F_{p}\to F(t)$ (Section~\ref{sec:reference-run}).
\end{enumerate}

\section{Validity, accuracy and practical guidance}
\label{sec:limitations}

\subsection{The overlap criterion}

The Hertz--Mindlin law is a linear-elastic, half-space solution.  Its derivation
assumes that the contact radius is small compared with the sphere radius,
$a\ll R$, which through~\eqref{eq:contact-radius} is equivalent to
\begin{equation}
  \boxed{\;
  \frac{\delta_{n}}{R} \ll 1,
  \qquad\text{in practice}\qquad
  \frac{\delta_{n}}{R}\lesssim 1\% \;}
  \label{eq:overlap-criterion}
\end{equation}
Combining~\eqref{eq:delta-static} with the fact that a load $F$ applied to a
platen of width $B$ is carried by roughly $n_{c}\approx B/(2\bar{R})$ contacts
gives an \emph{a priori} estimate of the overlap in the top layer,
\begin{equation}
  \frac{\delta_{n}}{\bar{R}}
  \approx \frac{1}{\bar{R}}
  \left(\frac{3F}{4 n_{c} E^{*}\sqrt{\bar{R}/2}}\right)^{2/3}.
  \label{eq:overlap-estimate}
\end{equation}
Table~\ref{tab:overlap} evaluates~\eqref{eq:overlap-estimate} for the reference
material.  It shows that with $E=\qty{5}{\mega\pascal}$ the model is
quantitatively trustworthy only up to a few tens of newtons; beyond that the
grains interpenetrate visibly and the computed stiffness of the packing is no
longer that of a Hertzian assembly.  This is a limitation of the chosen
\emph{parameters}, not of the implementation, and it is quantified further in
Section~\ref{sec:load-level}.

\begin{table}[htbp]
  \centering
  \caption{Estimated overlap in the layer directly under the platen, from
  Eq.~\eqref{eq:overlap-estimate}, for $B=\qty{0.2}{\metre}$,
  $\bar{R}=\qty{12}{\milli\metre}$, $n_{c}\approx8$, $E^{*}=\qty{2.67}{\mega\pascal}$.}
  \label{tab:overlap}
  \begin{tabular}{S[table-format=4.0]S[table-format=3.1]S[table-format=2.2]l}
    \toprule
    {$F$ [\unit{\newton}]} & {$F/n_{c}$ [\unit{\newton}]} &
    {$\delta_{n}/\bar{R}$ [\%]} & Assessment \\
    \midrule
      1  &  0.1 &  0.42 & quantitatively valid \\
     10  &  1.3 &  2.3  & acceptable \\
    100  & 12.5 & 10.6  & qualitative only \\
    1000 &  125 & 49.2  & outside the model's range \\
    \bottomrule
  \end{tabular}
\end{table}

\subsection{Artificial softening}

Real sand or glass has $E\sim\qty{50}{\giga\pascal}$; the reference material
uses $\qty{5}{\mega\pascal}$, four orders of magnitude softer.  This is the usual
compromise: from~\eqref{eq:rayleigh}, $\Delta t \propto E^{-1/2}$, so restoring a
realistic stiffness would divide the time step by about $100$ and multiply the
run time accordingly.  From~\eqref{eq:delta-static},
$\delta_{n}\propto E^{-2/3}$, so the softening inflates the overlaps by a factor
of about $460$ at a given load.  When a quantitative stress--strain response is
required, the modulus must be raised until the results become insensitive to it,
which is the softening-convergence study that any serious application should
include.

\subsection{Other sources of error}

\begin{itemize}[leftmargin=1.6em]
  \item \emph{Time-step error}: halving $\Delta t$ should change the packing
        statistics by less than the run-to-run scatter due to the random seed.
        Because a granular assembly is chaotic, individual trajectories will
        never converge; only ensemble quantities (porosity, coordination number,
        load--displacement curve) do.
  \item \emph{Tangential history rotation}: the projection
        \eqref{eq:xi-project} does not conserve the magnitude of
        $\boldsymbol{\xi}$ when the contact plane rotates quickly; the error is
        $O(\Delta t^{2})$ per step but accumulates in long-lived rolling
        contacts.
  \item \emph{Absence of rolling resistance}: a two-dimensional assembly of
        smooth discs is unrealistically free to roll, which under-predicts the
        shear strength and the angle of repose.  Adding a rolling-friction
        torque is the single most valuable extension of the present model.
  \item \emph{Boundary effects}: with $B/\bar{d}\approx 8$ grain diameters
        across, wall friction carries a significant share of the applied load
        (Janssen effect), and the measured platen reaction is not a pure
        material property.  A wider box or periodic lateral boundaries are
        required for constitutive characterisation.
  \item \emph{Two dimensions}: force chains in 2D are more percolating and the
        packing fraction is systematically higher than in 3D; results transfer
        qualitatively, not quantitatively.
\end{itemize}
